\section{Clipped Importance Sampling Policy Optimization}

Clipped Importance Sampling Policy Optimization (CISPO~\citep{minimax2025minimaxm1scalingtesttimecompute}) propsed in MiniMax-M1 optimizes a policy-gradient objective using a clipped importance sampling ratio as a stop-gradient weight.

The CISPO objective is defined as
\[
\mathcal{J}_{\mathrm{CISPO}}(\theta)
=
\mathbb{E}_{(q,a)\sim \mathcal{D},\ \{o_i\}_{i=1}^{G}\sim \pi_{\theta_{\mathrm{old}}}(\cdot\mid q)}
\left[
    \frac{1}{\sum_{i=1}^{G}|o_i|}
    \sum_{i=1}^{G}
    \sum_{t=1}^{|o_i|}
    \operatorname{sg}\!\left(\hat r_{i,t}(\theta)\right)
    \hat A_{i,t}
    \log \pi_\theta(o_{i,t}\mid q,o_{i,<t})
\right].
\]

Here, $\hat A_{i,t}$ is the estimated advantage for token $o_{i,t}$, and $\operatorname{sg}(\cdot)$ denotes the stop-gradient operation. Therefore, $\operatorname{sg}\!\left(\hat r_{i,t}(\theta)\right)$ is treated as a constant when taking gradients with respect to $\theta$. The clipped importance sampling ratio is
\[
    \hat r_{i,t}(\theta)
    =
    \operatorname{clip}
    \left(
        r_{i,t}(\theta),
        1-\epsilon_{\mathrm{low}},
        1+\epsilon_{\mathrm{high}}
    \right),
\]
where the token-level importance sampling ratio is
\[
    r_{i,t}(\theta)
    =
    \frac{
        \pi_\theta(o_{i,t}\mid q,o_{i,<t})
    }{
        \pi_{\theta_{\mathrm{old}}}(o_{i,t}\mid q,o_{i,<t})
    }.
\]

Equivalently,
\[
    \hat r_{i,t}(\theta)
    =
    \begin{cases}
    1-\epsilon_{\mathrm{low}},
    &
    r_{i,t}(\theta) < 1-\epsilon_{\mathrm{low}},
    \\[4pt]
    r_{i,t}(\theta),
    &
    1-\epsilon_{\mathrm{low}}
    \leq r_{i,t}(\theta)
    \leq 1+\epsilon_{\mathrm{high}},
    \\[4pt]
    1+\epsilon_{\mathrm{high}},
    &
    r_{i,t}(\theta) > 1+\epsilon_{\mathrm{high}}.
    \end{cases}
\]

Thus, CISPO does not clip the surrogate objective itself. Instead, it clips the token-level importance sampling ratio and uses the clipped ratio as a stop-gradient coefficient in front of the log-probability term.

\subsection{Gradient Difference between PPO and CISPO}

The main difference between PPO and CISPO is how clipping affects the policy gradient.

For PPO, the clipped surrogate objective at token $(i,t)$ can be written as
\[
    L^{\mathrm{PPO}}_{i,t}(\theta)
    =
    \min
    \left(
        r_{i,t}(\theta)\hat A_{i,t},
        \operatorname{clip}
        \left(
            r_{i,t}(\theta),
            1-\epsilon,
            1+\epsilon
        \right)
        \hat A_{i,t}
    \right),
\]
where
\[
    r_{i,t}(\theta)
    =
    \frac{
        \pi_\theta(o_{i,t}\mid q,o_{i,<t})
    }{
        \pi_{\theta_{\mathrm{old}}}(o_{i,t}\mid q,o_{i,<t})
    }.
\]

Consider the case
\[
    \hat A_{i,t} > 0
    \quad \text{and} \quad
    r_{i,t}(\theta) > 1+\epsilon.
\]
Then the clipped PPO objective becomes
\[
    L^{\mathrm{PPO}}_{i,t}(\theta)
    =
    (1+\epsilon)\hat A_{i,t}.
\]
This term is constant with respect to $\theta$, and therefore
\[
    \nabla_\theta L^{\mathrm{PPO}}_{i,t}(\theta)
    =
    0.
\]
Thus, when the probability of a positive-advantage token has already increased beyond the clipping range, PPO removes the gradient contribution from this token.

Similarly, if
\[
    \hat A_{i,t} < 0
    \quad \text{and} \quad
    r_{i,t}(\theta) < 1-\epsilon,
\]
then the clipped PPO objective becomes
\[
    L^{\mathrm{PPO}}_{i,t}(\theta)
    =
    (1-\epsilon)\hat A_{i,t},
\]
which is also constant with respect to $\theta$. Hence,
\[
    \nabla_\theta L^{\mathrm{PPO}}_{i,t}(\theta)
    =
    0.
\]
In this case, PPO removes the gradient contribution when the probability of a negative-advantage token has already decreased beyond the clipping range.

By contrast, CISPO uses the clipped importance sampling ratio as a stop-gradient coefficient:
\[
    \mathcal{J}_{\mathrm{CISPO}}(\theta)
    =
    \mathbb{E}
    \left[
        \frac{1}{\sum_{i=1}^{G}|o_i|}
        \sum_{i=1}^{G}
        \sum_{t=1}^{|o_i|}
        \operatorname{sg}\!\left(\hat r_{i,t}(\theta)\right)
        \hat A_{i,t}
        \log \pi_\theta(o_{i,t}\mid q,o_{i,<t})
    \right].
\]

Since
\[
    \operatorname{sg}\!\left(\hat r_{i,t}(\theta)\right)
\]
is treated as a constant during backpropagation, the gradient of the CISPO objective is
\[
    \nabla_\theta \mathcal{J}_{\mathrm{CISPO}}(\theta)
    =
    \mathbb{E}
    \left[
        \frac{1}{\sum_{i=1}^{G}|o_i|}
        \sum_{i=1}^{G}
        \sum_{t=1}^{|o_i|}
        \operatorname{sg}\!\left(\hat r_{i,t}(\theta)\right)
        \hat A_{i,t}
        \nabla_\theta
        \log \pi_\theta(o_{i,t}\mid q,o_{i,<t})
    \right].
\]

Therefore, even when the ratio $r_{i,t}(\theta)$ falls outside the clipping interval, CISPO does not make the gradient of that token vanish. Instead, the clipped ratio $\hat r_{i,t}(\theta)$
controls the magnitude of the gradient, while the token still contributes a policy-gradient signal through
\[
    \nabla_\theta \log \pi_\theta(o_{i,t}\mid q,o_{i,<t}).
\]

In short, PPO clips the surrogate objective, which can cause clipped tokens to have zero gradient. CISPO clips the importance sampling ratio and uses it as a stop-gradient weight, so clipped tokens can still contribute to the policy gradient.